\chapter{Návrh aplikace}
\label{ch:navrh}

% TODO: úvodní odstavec, který stručně nastíní nadcházející kapitoly v Návrhu aplikace,
% což bude Existující řešení, Uživatelská studie (strukturované rozhovory),
% z nich vyplynou Use cases, FR/NFR, případy užití,
% a na jejichž základě bude zkonstruován HiFi prototyp aplikace.

\section{Existující řešení}
\label{sec:existujici}

Mobilní aplikace pro sledování stravy a tělesné hmotnosti patří v~současnosti k~nejčastěji používaným nástrojům v~oblasti digitálního zdraví.
Přehledové studie ukazují, že nutriční a dietní aplikace jsou běžnou součástí změny životního stylu a mohou přispět k~lepší kontrole stravovacích návyků a k~mírné redukci hmotnosti~\cite{nogueira2024mobileApplications}.

Současné analýzy trhu s~digitálními nástroji pro výživu potvrzují, že v~hlavních mobilních obchodech existují desítky až stovky aplikací zaměřených na sledování příjmu potravy, plánování jídelníčku a takzvanou precision nutrition.
Abeltino a~kol. ve svém přehledu shrnují, že pro běžné uživatele i odborníky je k~dispozici široké spektrum aplikací pro monitorování stravy a tvorbu jídelních plánů~\cite{abeltino2025digitalApplications}.
Samad a~kol. při systematickém průzkumu tří velkých obchodů s~aplikacemi identifikovali 473 aplikací souvisejících se sledováním příjmu potravy a doporučováním jídel, z~nichž 80 podrobně analyzovali~\cite{samad2022smartphoneApps}.
To dokládá, že i v~jedné poměrně úzce vymezené kategorii existují desítky až stovky vzájemně konkurujících aplikací.

Přehledy výživových aplikací ukazují, že většina z~nich staví na velmi podobném jádru a liší se spíše v~uživatelském rozhraní a doplňkových funkcích než v~samotné logice sledování stravy~\cite{nogueira2024mobileApplications,abeltino2025digitalApplications}.
Typická aplikace pro sledování stravy nabízí zejména:

\begin{itemize}
	\item potravinový deník pro každodenní záznam jídel,
	\item databázi potravin s~nutričními hodnotami a fulltextovým vyhledáváním,
	\item možnost ukládat vlastní recepty a oblíbená jídla,
	\item nastavení cílového energetického příjmu a sledování makroživin,
	\item základní zpětnou vazbu ve formě grafů a statistik~\cite{gioia2023mobileApps}.
\end{itemize}

Z~hlediska interakce s~uživatelem je klíčové zadávání dat.
Většina komerčních aplikací stále vyžaduje manuální zapisování jídel prostřednictvím vyhledávání v~databázi nebo zadávání receptů.
Částečnou automatizaci přinášejí:

\begin{itemize}
	\item skenování čárových kódů balených potravin,
	\item předvyplněné šablony jídel a možnost kopírovat dříve zapsaná jídla,
	\item notifikace připomínající záznam jednotlivých jídel~\cite{abeltino2025digitalApplications}.
\end{itemize}

Systematické studie zároveň upozorňují, že ruční zapisování je časově náročné a v~praxi často vede k~nekompletním či zkresleným záznamům.
Li a~kol. ukazují, že výstupy nutričních aplikací se mohou od referenčních metod významně lišit, a to jak u~manuálního zapisování, tak u~řešení s~podporou rozpoznání obrazu.
Energetický příjem bývá v~závislosti na typu stravy nadhodnocen nebo podhodnocen o~stovky kilodžaulů až kilokalorií denně~\cite{li2024evaluating}.

V~reakci na omezení manuálního sběru dat se začíná prosazovat využití umělé inteligence a počítačového vidění.
Samad a~kol. ukazují, že pouze malá část komerčních aplikací nabízí pokročilé funkce automatického rozpoznávání jídel z~fotografie a odhadu porce, a že v~souboru 80 hodnocených aplikací byla jen jedna, která dokázala z~fotografie automaticky určit druh jídla i velikost porce~\cite{samad2022smartphoneApps}.
Aburub a~kol. ve svém přehledu food scanning aplikací shrnují, že skenování čárových kódů je dnes běžnou funkcí, zatímco plně automatické rozpoznávání talíře zůstává výjimkou a typicky vyžaduje potvrzení a korekci ze strany uživatele~\cite{aburub2025foodScanning}.

Nogueira-Rio a~kol. zdůrazňují, že současná generace výživových aplikací stojí na kombinaci tradičního self monitoringu a nových AI nástrojů, přičemž zásadní výzvou je skloubit vyšší míru automatizace s~přijatelnou přesností a použitelností v~každodenním životě~\cite{nogueira2024mobileApplications}.

Pro tuto práci jsou reprezentativní zejména dvě aplikace, které pokrývají globální i lokální kontext.
První je mezinárodní aplikace MyFitnessPal a druhou jsou československé Kalorické Tabulky (Dine4Fit).
Následující podkapitoly je stručně představují a vymezují jejich silné a slabé stránky ve vztahu k~cíli této diplomové práce.

\subsection{MyFitnessPal}
\label{subsec:myfitnesspal}

MyFitnessPal patří mezi nejznámější a nejrozšířenější aplikace pro sledování stravy a tělesné hmotnosti na světě.
Je dostupná pro Android, iOS i webové rozhraní a dlouhodobě se pohybuje na předních příčkách kategorií Health \& Fitness v~hlavních mobilních obchodech~\cite{myFitnessPal2025}.

\subsubsection*{Hlavní funkcionalita}

Základem MyFitnessPal je potravinový deník navázaný na rozsáhlou databázi potravin.
Oficiální materiály uvádějí databázi s~miliony položek včetně běžných potravin, značkových výrobků i restaurací.
Uživatel může zapisovat jídelníček několika způsoby:

\begin{itemize}
	\item vyhledáváním v~databázi podle názvu potraviny,
	\item výběrem z~historie a často používaných jídel,
	\item skenováním čárového kódu u~balených potravin,
	\item zadáváním vlastních receptů a pokrmů.
\end{itemize}

Aplikace průběžně počítá celkový denní energetický příjem a zobrazuje rozložení makroživin, jako jsou sacharidy, tuky a bílkoviny.
V~placené verzi umožňuje sledovat i další nutriční parametry, například vlákninu, cukr, sůl nebo vybrané mikronutrienty.
Součástí systému je také záznam fyzických aktivit, sledování tělesné hmotnosti a integrace s~externími zařízeními a službami, jako jsou Apple Health nebo fitness náramky~\cite{myFitnessPal2025}.

Obchodní model aplikace odpovídá freemium přístupu.
Základní funkce potravinového deníku a databáze potravin jsou dostupné zdarma, rozšířené analytické nástroje, pokročilé nastavení cílů, detailnější nutriční přehledy a některé zrychlené způsoby zadávání jídel jsou součástí placeného předplatného MyFitnessPal Premium~\cite{myFitnessPal2025}.

\subsubsection*{Rozpoznávání jídel z~fotografie}

MyFitnessPal je jedním z~prvních velkých komerčních produktů, které nasadily ve větším měřítku foto rozpoznávání jídel.
Funkce Meal Scan umožňuje uživateli namířit fotoaparát na talíř a aplikace na základě obrazu navrhne odpovídající položky z~databáze.
Podle oficiální dokumentace jde o~funkci dostupnou v~rámci předplatného Premium a aktuálně podporovanou pro vybrané verze iOS a Androidu~\cite{myFitnessPal2025}.

Výrobce uvádí, že Meal Scan využívá metody počítačového vidění a strojového učení k~identifikaci více položek na jednom talíři a jejich propojení s~interní databází.
V~praxi funguje Meal Scan jako asistovaná funkce: uživatel dostane seznam odpovídajících jídel, zvolí nejbližší položku a případně upraví velikost porce~\cite{myFitnessPal2025}.

\subsubsection*{Silné stránky a omezení}

Z~hlediska této práce jsou klíčové následující výhody MyFitnessPal:

\begin{itemize}
	\item rozsáhlá databáze potravin pokrývající globální trh, včetně restaurací a značkových výrobků~\cite{myFitnessPal2025},
	\item kombinace více vstupních metod (vyhledávání, historie, čárový kód, foto rozpoznávání)~\cite{myFitnessPal2025},
	\item dlouhodobá přítomnost na trhu, díky níž existuje řada nezávislých hodnocení této kategorie nástrojů, včetně studií zaměřených na kvalitu a funkce aplikací pro sledování stravy~\cite{samad2022smartphoneApps}.
\end{itemize}

Současná literatura i uživatelské zkušenosti ale upozorňují na několik omezení:

\begin{itemize}
	\item část databáze tvoří uživateli přidané položky, což může vést k~chybám v~nutričních údajích a nutnosti ruční korekce, například po skenování čárového kódu~\cite{li2024evaluating},
	\item i přes použití foto rozpoznávání zůstává nutnost ruční kontroly a úprav, zejména u~složitějších jídel a kombinovaných receptů~\cite{li2024evaluating},
	\item aplikace je navržena primárně pro globální trh, a proto nemusí dostatečně pokrývat některé lokální produkty nebo značky, zejména ve středoevropském kontextu.
\end{itemize}

Pro návrh vlastního řešení je MyFitnessPal důležitým referenčním příkladem robustního ekosystému, který kombinuje manuální self monitoring s~prvky umělé inteligence.
Současně ukazuje, že i jedna z~nejrozšířenějších aplikací využívá umělou inteligenci primárně jako doplněk, nikoli jako plnou náhradu manuálního zadávání.

\subsection{Kalorické Tabulky}
\label{subsec:kalorickeTabulky}

Kalorické Tabulky, vyvíjené společností Dine4Fit~a.s., představují nejpoužívanější aplikaci pro sledování stravy v~českém a slovenském prostředí.
Podle oficiálních informací firmy se stala nejpopulárnější aplikací na hlídání stravy v~Česku i na Slovensku a měsíčně ji používá více než jeden milion uživatelů~\cite{dine4fit2025}.
Tyto údaje potvrzuje také nezávislá reportáž na serveru CzechCrunch, která uvádí více než 1,5~milionu aktivních uživatelů a několik milionů registrací~\cite{houska2024kalorickeTabulky}.

Aplikace je dostupná na webu i jako mobilní aplikace pro Android a iOS a podporuje také vybrané chytré hodinky, například Garmin nebo Apple Watch~\cite{dine4fit2025}.

\subsubsection*{Hlavní funkcionalita}

Základní princip Kalorických Tabulek se podobá globálním aplikacím pro sledování stravy, je však výrazně přizpůsoben specifikům českého a slovenského trhu.
Klíčové funkce jsou:

\begin{itemize}
	\item databáze potravin zaměřená na produkty dostupné v~Česku a na Slovensku, včetně značkových výrobků, restaurací a uživatelských receptů~\cite{dine4fit2025},
	\item potravinový deník s~možností zapisovat jednotlivá jídla v~průběhu dne a sledovat energetický příjem v~kcal nebo kJ,
	\item přehled denního příjmu a výdeje energie, sledování pitného režimu a tělesné hmotnosti~\cite{dine4fit2025},
	\item podpora záznamu fyzických aktivit a jejich vlivu na energetickou bilanci.
\end{itemize}

Významnou součástí ekosystému je i kalorická kalkulačka, která odvozuje doporučený energetický příjem z~tělesných parametrů a cílů uživatele.
Dine4Fit uvádí, že nastavení poměru makroživin vychází z~odborných doporučení a umožňuje zvolit různé cíle, jako je udržení hmotnosti, hubnutí nebo nárůst svalové hmoty~\cite{dine4fit2025}.

Aplikace používá freemium model.
Základní funkcionalita, tedy deník, výpočet energie a základní nutriční hodnoty, je dostupná zdarma, zatímco detailnější analýzy, rozšířené sledování mikronutrientů, speciální jídelníčky a pokročilé statistiky jsou součástí placeného předplatného Premium~\cite{dine4fit2025,houska2024kalorickeTabulky}.

\subsubsection*{Rozpoznávání jídel z~fotografie}

Novější verze Kalorických Tabulek integrují funkci rozpoznávání jídel z~fotografie pomocí umělé inteligence.
Podle popisu aplikace v~Google Play může uživatel jídlo vyfotit a aplikace navrhne odpovídající potraviny a jejich přibližnou gramáž, čímž automaticky spočítá kalorie a živiny~\cite{dine4fit2025}.
Dine4Fit tuto funkci dále popisuje na svém webu a blogu jako prémiovou novinku využívající umělou inteligenci pro rozpoznání jídel a odhad hmotnosti na základě fotografie~\cite{dine4fit2025}.

Technologické detaily implementace nejsou veřejně detailně popsány, mediální články nicméně zmiňují, že aplikace využívá detekci potravin pomocí umělé inteligence ve velkém objemu zaznamenaných jídel a že funkce je nasazována postupně s~důrazem na praktické využití v~běžném hubnutí~\cite{houska2024kalorickeTabulky}.

Stejně jako u~MyFitnessPal nejde o~plně autonomní systém.
Uživatelé musí výsledek zkontrolovat, případně upravit typ jídla nebo velikost porce.
Odborná literatura zatím neobsahuje nezávislou kvantitativní studii přesnosti této konkrétní funkce, což odpovídá obecnějšímu trendu, kdy detailní vědecká validace komerčních AI funkcí obvykle následuje s~časovým odstupem za jejich praktickým nasazením~\cite{samad2022smartphoneApps,li2024evaluating,aburub2025foodScanning}.

\subsubsection*{Silné stránky a omezení}

V~českém a slovenském kontextu mají Kalorické Tabulky několik výrazných výhod:

\begin{itemize}
	\item silná lokalizace databáze potravin na regionální produkty, včetně běžných značek a lokálních jídel~\cite{dine4fit2025},
	\item dostupnost webového rozhraní i mobilních aplikací, což usnadňuje zapisování v~různých situacích~\cite{dine4fit2025},
	\item velká uživatelská základna, která přispívá k~rozšiřování databáze a komunitním receptům~\cite{houska2024kalorickeTabulky}.
\end{itemize}

Současně však z~uživatelských recenzí a odborných diskusí vyplývá několik slabších míst relevantních pro návrh této práce:

\begin{itemize}
	\item práce se složitějšími recepty může být časově náročná a vyžaduje manuální sestavování položek,
	\item podobně jako u~jiných komunitně rozšiřovaných databází se mohou vyskytnout chyby v~nutričních údajích, pokud uživatelé zadávají data nepřesně~\cite{li2024evaluating},
	\item funkce rozpoznávání jídel pomocí umělé inteligence je relativně nová a dostupné informace ji popisují jako doplněk k~tradičnímu zadávání, nikoli jako plnou náhradu manuálního zápisu~\cite{dine4fit2025,houska2024kalorickeTabulky}.
\end{itemize}

Pro navrhovanou aplikaci jsou Kalorické Tabulky cenným vzorem z~hlediska lokalizace, práce s~regionálně specifickou databází a integrace funkcí umělé inteligence do jinak klasického stravovacího deníku.

\subsection{Shrnutí aplikací na trhu}
\label{subsec:shrnutiTrhu}

Odborná literatura a analýzy komerčního trhu ukazují, že výživové a dietní aplikace tvoří velmi početnou a relativně homogenní kategorii digitálních nástrojů.
Přehledy trhu a analýzy obchodů s~aplikacemi naznačují, že v~globálním měřítku existují desítky až stovky aplikací zaměřených na sledování stravy, které si vzájemně konkurují v~rámci několika málo funkčních podkategorií, například calorie tracker, nutrition tracker, diet planner nebo food scanning aplikace~\cite{abeltino2025digitalApplications,samad2022smartphoneApps}.

Samad a~kol. identifikovali 473 aplikací souvisejících s~příjmem potravy a doporučováním jídel a 80 z~nich podrobili detailnímu hodnocení~\cite{samad2022smartphoneApps}.
Gioia a~kol. pak ve svém přehledu vybrali reprezentativní soubor dietních aplikací a analyzovali jejich funkcionalitu a vhodnost pro výzkum stravovacích návyků~\cite{gioia2023mobileApps}.
Přestože se metodiky liší, obě práce ukazují, že ve sledované kategorii se vyskytují desítky až stovky aplikací, které sdílejí velmi podobnou základní funkcionalitu.

Napříč těmito aplikacemi se opakují zejména tyto prvky:

\begin{itemize}
	\item potravinový deník pro každodenní záznam jídel,
	\item databáze potravin s~nutričními hodnotami,
	\item ruční nebo poloautomatické zadávání pomocí čárového kódu a šablon,
	\item přehledy o~denním příjmu energie a makroživin, někdy doplněné o~motivační prvky a připomenutí,
	\item napojení na sledování tělesné hmotnosti a často i na fyzickou aktivitu~\cite{abeltino2025digitalApplications,gioia2023mobileApps}.
\end{itemize}

Rozdíly mezi jednotlivými aplikacemi se proto často týkají spíše míry propracovanosti uživatelského rozhraní a motivace, kvality a lokalizace potravinové databáze, míry integrace s~nositelnou elektronikou a rozsahu využití umělé inteligence pro automatizaci záznamu~\cite{nogueira2024mobileApplications,abeltino2025digitalApplications}.

Studie zaměřené na přesnost a uživatelskou zátěž potvrzují, že ruční zapisování stravy je pro mnoho lidí dlouhodobě obtížné a že dietní aplikace mohou energetický příjem podhodnocovat nebo nadhodnocovat, někdy i o~stovky kilokalorií denně~\cite{li2024evaluating}.
To vytváří prostor pro inovace v~oblasti automatizace sběru dat, především prostřednictvím obrazového rozpoznávání a dalších metod umělé inteligence.
Současně však scoping a systematické přehledy upozorňují, že plně automatické foto rozpoznávání je zatím spíše výjimkou a že i pokročilé funkce vyžadují aktivní roli uživatele~\cite{samad2022smartphoneApps,aburub2025foodScanning}.

MyFitnessPal a Kalorické Tabulky představují dva výrazné příklady této situace.
Na jedné straně ukazují vyspělost trhu, tedy rozsáhlé databáze, škálu vstupních metod, robustní ekosystém a milionové uživatelské základny~\cite{myFitnessPal2025,dine4fit2025,houska2024kalorickeTabulky}.
Na straně druhé potvrzují, že základní model zůstává obdobný jako v~počátcích mobilních dietních aplikací.
Uživatel stále primárně ručně zapisuje jídla do deníku a funkce foto rozpoznávání slouží jako doplněk, nikoli jako plná náhrada manuálního sběru dat~\cite{li2024evaluating,aburub2025foodScanning}.

Pro návrh aplikace v~této diplomové práci z~toho vyplývají tři hlavní závěry:

\begin{enumerate}
	\item Trh již nabízí velké množství aplikací se srovnatelnou základní funkcionalitou, takže inovace by měla cílit na konkrétní slabá místa, například uživatelskou zátěž při zadávání nebo kvalitu a lokalizaci databáze.
	\item Zvýšení míry automatizace pomocí počítačového vidění a umělé inteligence je slibnou cestou, ale vyžaduje pečlivé řešení otázky přesnosti a transparentní komunikaci nejistot vůči uživateli.
	\item Lokální kontext a kvalita potravinové databáze mohou být v~českém prostředí stejně důležité jako samotná technická vyspělost modelu rozpoznávání.
\end{enumerate}

\section{Uživatelská studie}
\label{sec:uzivatelskaStudie}

Uživatelská studie představuje soubor metod, jejichž cílem je porozumět potřebám, motivacím a kontextu cílových uživatelů a tyto poznatky následně převést do návrhu systému.
V~oblasti mobilního zdraví je uživatelsky orientovaný přístup považován za důležitý, protože ovlivňuje použitelnost, důvěru, dlouhodobou adherenci a tím i reálný dopad aplikace~\cite{cornet2020untoldStories}.

V~této práci byla uživatelská studie použita jako hlavní vstup pro návrh aplikace pro rozpoznávání potravin pomocí umělé inteligence a pro sledování kalorického příjmu.
Cílem bylo identifikovat situace, ve kterých je ruční zapisování pro uživatele překážkou, jaká míra nepřesnosti je akceptovatelná, jak uživatelé vnímají selhání umělé inteligence a jaké podpůrné funkce zvyšují pravděpodobnost dlouhodobého používání.

\subsection{Metodika}
\label{subsec:metodika}

Jako metoda sběru dat byly zvoleny strukturované rozhovory s~uživateli z~cílové skupiny.
Rozhovory byly vedeny tak, aby pokryly současnou praxi zapisování jídel, očekávání od rozpoznávání pomocí umělé inteligence a také preference v~oblasti přehledů, motivace, práce s~daty a integrací.
Tento postup odpovídá doporučením v~literatuře, kde jsou kvalitativní techniky typu rozhovorů nebo skupinových diskusí uváděny jako vhodný zdroj požadavků pro návrh aplikací v~oblasti zdraví~\cite{molina2020proposalUserCentered}.

Rozhovory byly zpracovány do samostatných reportů a dále analyzovány tematickým způsobem.
V~této kapitole jsou výsledky shrnuty do hlavních témat a následně formalizovány do scénářů případů užití a do sady funkčních a nefunkčních požadavků.
Reporty v~plném znění jsou součástí příloh této práce.

\subsection{Charakteristika participantů}
\label{subsec:participanti}

Pro rozhovory byli vybráni čtyři participanti, kteří reprezentují rozdílné cíle a kontext používání.
V~textu jsou uváděni anonymizovaně jako P1 až P4.

\begin{description}
	\item[P1] muž ve věku 27~let, jehož cílem je nabírání svalové hmoty.
		Cvičí přibližně třikrát týdně a preferuje sledování makroživin.
		V~minulosti aplikaci pro zapisování využíval, ale přestal, když dosáhl cíle.
		Je ochoten tolerovat určitou míru nepřesnosti, pokud mu aplikace ušetří práci.
	\item[P2] muž ve věku 24~let, intenzivně cvičí a cíleně udržuje kalorický nadbytek.
		U~zapisování klade důraz na přesnost a důvěryhodnost, preferuje rychlé metody zadávání u~balených potravin, typicky čárové kódy.
		U~jídel z~restaurace akceptuje typicky vyšší nepřesnost při zápisu než u~jiných pokrmů.
	\item[P3] žena ve věku 54~let, dlouhodobě se snaží zhubnout a zohledňuje zdravotní souvislosti.
		Zkoušela existující aplikace, ale odradila ji časová náročnost, časté chybění položek v~databázi a nutnost odhadovat gramáž, zejména u~domácího vaření.
	\item[P4] žena ve věku 28~let, která dříve své jídlo pravidelně vážila, nyní častěji zapisuje odhadem a upřednostňuje jednoduché předvolby porcí.
		Má jasnou preferenci, aby fotografie byla primárním vstupem, ale současně požaduje kvalitní záložní cestu při selhání.
		Zásadní je pro ni transparentnost umělé inteligence a rozlišení běžných limitů modelu od technické chyby aplikace.
\end{description}

\subsection{Výsledky rozhovorů}
\label{subsec:vysledkyRozhovoru}

\paragraph{Rychlost zápisu a dlouhodobá udržitelnost.}
Opakovaným tématem napříč rozhovory je, že dlouhodobé používání je přímo podmíněno rychlostí a jednoduchostí.
P1 zmiňuje cílový stav rychlého zápisu jednoduché položky přibližně do šesti kliknutí bez dodatečné editace.
P2 uvádí podobnou hranici pro rychlé položky.
P4 toleruje pomalejší zpracování zejména na začátku, ale očekává postupné zrychlení a udržení zápisu v~přiměřeném čase.
Zcela klíčové je to pro P3, která z~důvodů pomalosti procesu zápisu pokaždé přestala aplikaci používat, a to ještě před dosažením svých zdravotních cílů.

\paragraph{Akceptovatelná nepřesnost a důvěra v~umělou inteligenci.}
Participanti se liší v~toleranci nepřesnosti, což je klíčové pro definici cílových metrik přesnosti.
P1 je ochoten tolerovat přibližnou odchylku do dvaceti procent, pokud mu aplikace výrazně zjednoduší proces.
Naopak P2 požaduje nepřesnost pouze v~řádu jednotek procent a bez důvěryhodného výsledku považuje funkci za málo užitečnou.
P3 chápe omezení vysoké přesnosti, nicméně požaduje, aby výstup byl srozumitelný a kontrolovatelný, zejména u~domácích jídel.

\paragraph{Selhání rozpoznání.}
Pro praktickou použitelnost funkce rozpoznávání je kritické eliminovat situace, kdy se uživatel ocitne ve slepé uličce.
P4 požaduje, aby při selhání fotografie bylo možné plynule přejít na textový popis a aby bylo možné rozpoznání znovu vyvolat z~rozpracovaného záznamu bez zakládání nového záznamu.
P1 zároveň navrhuje možnost doplnit k~fotografii krátkou doplňující informaci pro lepší výsledek.
Důležitým aspektem je také rozlišení, zda jde o~limit rozpoznání, nebo o~technickou chybu aplikace, a nabídnutí vhodného dalšího kroku.

\paragraph{Oblíbená jídla, historie a zkratky.}
Všichni participanti popisují situace, kdy jedí opakovaně podobná jídla, případně používají stejné ingredience.
P2 označuje oblíbené položky, historii a čtečku kódů za zásadní funkce.
P4 uvádí, že využívá oblíbené a duplikaci předešlého záznamu a ocení našeptávání dříve zadaných názvů.
P1 doplňuje potřebu rychlých zkratek, které umožní udržet zápis v~nízkém počtu kroků.

\paragraph{Domácí vaření a odhad množství.}
Domácí vaření se ukazuje jako jeden z~možných důvodů, proč uživatelé zapisování opouštějí.
P3 popisuje, že vaření od oka ztěžuje odhad gramáže a že největší bariéra nastává, když položka chybí v~databázi a je nutné ji vytvořit manuálně.
P2 zdůrazňuje potřebu práce s~plánovanými ingrediencemi a následnou jednoduchou úpravou skutečně použitých množství, včetně drobných ingrediencí typu olej nebo koření.
To vede k~požadavku na funkcionalitu, která podporuje plánování a následné rychlé dorovnání skutečně použitých množství surovin.

\paragraph{Kategorie jídel a práce se zpětným zápisem.}
Participanti se liší v~tom, zda zapisují před jídlem nebo zpětně.
P2 u~jídel z~restaurace preferuje doplnění až doma mimo sociální situaci.
P3 uvádí, že zapisovala i večer, případně plánovala dopředu, ale nedodržení plánů vedlo ke zdvojení práce.
Z~toho vyplývá potřeba podporovat jak okamžitý zápis, tak zpětný zápis z~galerie, a zároveň nevnucovat automatické zařazování kategorií podle času zadání.

\paragraph{Přehledy, motivace a dotazování nad daty.}
Základem je pro všechny denní přehled, tedy rychlá kontrola kalorií a makroživin.
P1 požaduje jednoduché přehledy po dnech, týdnech a měsících a možnost pokládat dotazy typu porovnání období podle makroživin.
P4 vnímá týdenní souhrn jako užitečné doplnění a navrhuje dotazování přirozeným jazykem s~personalizovanými doporučeními.
P2 uvádí, že motivační měsíční souhrny formou notifikace nebo e-mailu by mu pomohly udržet rutinu.
P3 klade důraz na dietní omezení a schopnost odpovědět, kolikrát došlo k~porušení diety s~proklikem na konkrétní dny.

\paragraph{Soukromí, ukládání dat a kontrola.}
Otázka soukromí se objevuje ve smyslu kontroly nad daty a transparentností práce s~nimi.
P1 požaduje jasné funkce pro mazání a informace o~tom, co se s~daty děje.
P4 je ochotna akceptovat automatické promazávání nejstarších záznamů mimo oblíbené položky.
P3 považuje za důležité mít možnost smazat jednotlivé záznamy, i když je pro ni umístění dat méně podstatné.

\paragraph{Integrace aktivity.}
Integrace se sportovními službami a hodinkami se opakuje u~všech participantů, byť s~rozdílnou hloubkou.
P1 chce jednoduché propojení a zobrazení příjmu versus výdeje.
P2 očekává přebírání spálených kalorií s~možností nastavit, zda a jak se výdej promítne do denního limitu.
P4 zmiňuje zájem o~integraci s~Garmin a zdůrazňuje, že ji zajímá zejména promítnutí spálených kalorií do denního přehledu.

\subsection{Scénáře případů užití}
\label{subsec:pripadyUziti}

V~návrhu aplikací se používají jak narativní popisy situací z~praxe, tak případy užití jako formální popis interakce aktéra se systémem.
V~této práci byly narativní situace využity při analýze rozhovorů, ale pro formální popis systému je dále uvedena pouze sada scénářů případů užití, aby nedocházelo k~duplicitě a aby byl výstup přímo použitelný pro návrh prototypu a testování.

\paragraph{UC01: Záznam jídla z~fotografie.}
\textbf{Charakteristika:} Uživatel chce rychle zapsat jídlo pomocí fotografie a následně výsledek zkontrolovat.
\textbf{Primární aktér:} uživatel.
Hlavní scénář:

\begin{enumerate}
	\item \emph{Uživatel} otevře funkci pro záznam jídla a pořídí fotografii.
	\item \emph{Systém} spustí rozpoznání, zobrazí návrh položek a odhad množství, případně indikaci nejistoty.
	\item \emph{Uživatel} zkontroluje návrh, upraví položky a množství, případně doplní chybějící položky.
	\item \emph{Uživatel} potvrdí uložení záznamu.
	\item \emph{Systém} uloží záznam a aktualizuje denní přehled.
\end{enumerate}

\paragraph{UC02: Oprava výsledku rozpoznání a opakování rozpoznání.}
\textbf{Charakteristika:} Uživatel chce opravit rozpoznané položky a v~případě potřeby vyvolat rozpoznání znovu bez zakládání nového záznamu.
\textbf{Primární aktér:} uživatel.
Hlavní scénář:

\begin{enumerate}
	\item \emph{Uživatel} otevře detail rozpoznaného jídla v~editační obrazovce.
	\item \emph{Systém} zobrazí seznam položek, množství, jednotky a souhrn živin.
	\item \emph{Uživatel} upraví názvy, množství, odstraní nebo přidá položky.
	\item \emph{Uživatel} zvolí možnost znovu rozpoznat, například po úpravě vstupu nebo doplnění informace.
	\item \emph{Systém} provede nové rozpoznání a aktualizuje návrh položek, zachová možnost ruční kontroly.
	\item \emph{Uživatel} potvrdí finální stav a uloží záznam.
\end{enumerate}

\paragraph{UC03: Ruční přidání záznamu jídla.}
\textbf{Charakteristika:} Uživatel chce vytvořit záznam bez použití fotografie nebo umělé inteligence.
\textbf{Primární aktér:} uživatel.
Hlavní scénář:

\begin{enumerate}
	\item \emph{Uživatel} otevře ruční zápis jídla.
	\item \emph{Uživatel} vyhledá položku v~databázi nebo zvolí položku z~historie.
	\item \emph{Uživatel} zadá množství a jednotku, gramy nebo kusy, případně předvolbu porce.
	\item \emph{Uživatel} potvrdí uložení záznamu.
	\item \emph{Systém} uloží záznam a aktualizuje denní přehled.
\end{enumerate}

\paragraph{UC04: Přidání balené potraviny pomocí čárového kódu.}
\textbf{Charakteristika:} Uživatel chce rychle přidat balenou potravinu načtením kódu.
\textbf{Primární aktér:} uživatel.
Hlavní scénář:

\begin{enumerate}
	\item \emph{Uživatel} otevře čtečku kódů a naskenuje čárový kód produktu.
	\item \emph{Systém} vyhledá produkt v~databázi a zobrazí nalezenou položku.
	\item \emph{Uživatel} zvolí množství a jednotku nebo předvolbu porce.
	\item \emph{Uživatel} potvrdí vložení do denního záznamu.
	\item \emph{Systém} uloží záznam a aktualizuje denní přehled.
\end{enumerate}

\paragraph{UC05: Přidání opakovaného jídla z~oblíbených nebo duplikací.}
\textbf{Charakteristika:} Uživatel chce co nejrychleji přidat jídlo, které se často opakuje.
\textbf{Primární aktér:} uživatel.
Hlavní scénář:

\begin{enumerate}
	\item \emph{Uživatel} otevře seznam oblíbených nebo historii záznamů.
	\item \emph{Systém} nabídne oblíbené položky, nedávné záznamy a našeptávání podle historie.
	\item \emph{Uživatel} vybere položku nebo duplikuje celý záznam jídla.
	\item \emph{Uživatel} případně upraví množství a potvrdí uložení.
	\item \emph{Systém} uloží nový záznam a aktualizuje denní přehled.
\end{enumerate}

\paragraph{UC06: Vaření s~režimem plán versus realita.}
\textbf{Charakteristika:} Uživatel chce při vaření zadat plánované ingredience a po dokončení upravit skutečně použité množství.
\textbf{Primární aktér:} uživatel.
Hlavní scénář:

\begin{enumerate}
	\item \emph{Uživatel} vytvoří záznam receptu nebo vaření a zadá plánované ingredience.
	\item \emph{Systém} průběžně počítá plánované souhrnné živiny receptu.
	\item \emph{Uživatel} po uvaření upraví skutečně použité množství, včetně drobných ingrediencí.
	\item \emph{Systém} přepočítá souhrny a nabídne uložení finální verze.
	\item \emph{Uživatel} uloží jídlo do denního záznamu a případně do vlastních receptů pro opakování.
\end{enumerate}

\paragraph{UC07: Nastavení profilu, cílů a dietních omezení.}
\textbf{Charakteristika:} Uživatel chce nastavit cíle, profil a kontext, který ovlivní interpretaci a doporučení.
\textbf{Primární aktér:} uživatel.
Hlavní scénář:

\begin{enumerate}
	\item \emph{Uživatel} otevře nastavení profilu.
	\item \emph{Uživatel} vyplní nebo upraví základní parametry a nastaví cílové hodnoty.
	\item \emph{Uživatel} nastaví dietní omezení nebo intolerance.
	\item \emph{Systém} uloží nastavení a promítne je do přehledů a upozornění.
\end{enumerate}

\paragraph{UC08: Přehledy, dotazy nad daty a export.}
\textbf{Charakteristika:} Uživatel chce vyhodnotit svůj progres a získat odpovědi nad historií, případně data sdílet.
\textbf{Primární aktér:} uživatel.
Hlavní scénář:

\begin{enumerate}
	\item \emph{Uživatel} otevře přehledy a zvolí období, týden nebo měsíc.
	\item \emph{Systém} zobrazí souhrny a trendy, případně upozorní na překročení limitů.
	\item \emph{Uživatel} položí dotaz nad daty, například porovnání období nebo dotaz na makroživiny.
	\item \emph{Systém} zobrazí odpověď a nabídne proklik na relevantní dny.
	\item \emph{Uživatel} volitelně spustí export dat do tabulkového formátu.
\end{enumerate}

\subsection{Funkční a nefunkční požadavky}
\label{subsec:pozadavky}

Z~rozhovorů a ze zadání práce byla sestavena sada funkčních a nefunkčních požadavků.
Požadavky jsou v~této kapitole uvedeny ve zkrácené podobě, aby byly použitelné jako přehledný podklad pro návrh prototypu a následnou implementaci.
Plné znění, včetně zdůvodnění a metody ověření, je uvedeno v~samostatném dokumentu specifikace, který je součástí příloh této práce.
Funkční požadavky na aplikaci jsou následující.

\paragraph{Základní evidence a denní přehled.}
\begin{description}[style=nextline]
	\item[FR-02] Denní přehled pro kalorie a makroživiny.
	\item[FR-03] Nastavení cílových hodnot, minimálně kalorie, volitelně makroživiny.
	\item[FR-04] Správa profilu uživatele.
	\item[FR-05] Kontrola nad daty a mazání jednotlivých záznamů, případně účtu.
\end{description}

\paragraph{Vstupy pro záznam jídla.}
\begin{description}[style=nextline]
	\item[FR-01] Ruční přidání záznamu jídla bez použití umělé inteligence.
	\item[FR-06] Pořízení fotografie jako vstup pro rozpoznání.
	\item[FR-12] Import fotografie z~galerie.
	\item[FR-14] Záznam jídla bez fotografie, například hlasem nebo textovým popisem.
\end{description}

\paragraph{Rozpoznávání pomocí umělé inteligence a práce s~nejistotou.}
\begin{description}[style=nextline]
	\item[FR-07] Návrh položek a porcí pomocí umělé inteligence.
	\item[FR-08] Indikace nejistoty výstupu.
	\item[FR-09] Vysvětlení limitů umělé inteligence a doporučení pro zlepšení vstupu.
	\item[FR-10] Rozlišení limitu umělé inteligence a technické chyby aplikace.
	\item[FR-11] Textový fallback při selhání rozpoznání.
	\item[FR-13] Opakování rozpoznání z~editační obrazovky.
\end{description}

\paragraph{Množství a jednotky.}
\begin{description}[style=nextline]
	\item[FR-15] Jednotky množství v~gramech i kusech, včetně předvoleb porcí.
\end{description}

\paragraph{Rychlé zkratky a práce s~historií.}
\begin{description}[style=nextline]
	\item[FR-16] Čtečka čárových kódů pro balené potraviny.
	\item[FR-17] Oblíbené položky.
	\item[FR-18] Duplikace předchozího záznamu.
	\item[FR-19] Našeptávání názvů podle historie.
\end{description}

\paragraph{Dietní kontext.}
\begin{description}[style=nextline]
	\item[FR-20] Dietní omezení a intolerance.
	\item[FR-21] Zobrazení porušení dietních omezení v~kalendáři.
\end{description}

\paragraph{Aktivita, přehledy a motivace.}
\begin{description}[style=nextline]
	\item[FR-22] Příjem versus výdej v~jednom pohledu.
	\item[FR-23] Nastavení integrace výdeje energie a jeho vlivu na limit.
	\item[FR-24] Týdenní a měsíční přehledy.
	\item[FR-27] Dotazy v~přirozeném jazyce nad daty.
	\item[FR-28] Měsíční motivační souhrn, notifikace nebo e-mail.
	\item[FR-29] Nastavitelné notifikace.
\end{description}

\paragraph{Data a provozní vlastnosti.}
\begin{description}[style=nextline]
	\item[FR-25] Export dat, například ve formátu CSV.
	\item[FR-26] Offline tolerance, koncept záznamu, dokončení později.
	\item[FR-30] Skrytí nebo zobrazení pokročilých funkcí v~uživatelském rozhraní.
\end{description}

\paragraph{Nefunkční požadavky.}
\begin{description}[style=nextline]
	\item[NFR-01] Akceptovatelná nepřesnost odhadu energie a makroživin nejvýše 10~\%.
	\item[NFR-02] Latence rozpoznání umělou inteligencí, výsledek nejvýše do 20~vteřin za běžných podmínek.
	\item[NFR-03] Počet kroků pro zápis jednoduché položky nejvýše 6 kroků bez dodatečné editace.
	\item[NFR-04] Časová náročnost zápisu, opakované jídlo do 1~minuty, nové jídlo do 5~minut.
	\item[NFR-05] Minimalistické a přehledné uživatelské rozhraní bez rušivých prvků.
\end{description}

\subsection{Shrnutí kapitoly}
\label{subsec:shrnutiUS}

Uživatelská studie ukázala, že hlavní bariérou dlouhodobého zapisování je kombinace časové náročnosti, nejistoty výsledků a slabé podpory pro opakované situace, zejména opakovaná jídla, balené produkty, domácí vaření a pokrmy z~restaurací.
Rozpoznávání pomocí umělé inteligence je vnímáno jako přínosné, ale pouze za předpokladu, že uživatel má kontrolu nad výsledkem a rozumí nejistotě spojené s~vyhodnocením.
Výstupy rozhovorů byly formulovány do scénářů případů užití a do sady požadavků, které tvoří přímý podklad pro návrh a realizaci prototypu.

\section{Prototyp aplikace}
\label{sec:prototyp}

Prototyp představuje zjednodušenou reprezentaci budoucího systému, která umožňuje ověřovat návrhová rozhodnutí dříve, než dojde k~nákladné implementaci.
V~praxi jde o~iterativní artefakt, jehož míra detailu se postupně zvyšuje a který slouží jak k~interní validaci návrhu, tak k~získání zpětné vazby od uživatelů~\cite{joyner2022lowFidelity,flohr2023contextBased}.

\subsection{Prototyp a související pojmy}
\label{subsec:pojmy}

Pojem prototyp v~kontextu návrhu uživatelského rozhraní zahrnuje různé formy reprezentace, od statických náčrtů a wireframů až po interaktivní prototypy.
Společným cílem je simulovat očekávané chování systému v~takovém rozsahu, aby bylo možné posoudit srozumitelnost toku obrazovek a kvalitu interakce~\cite{joyner2022lowFidelity,flohr2023contextBased}.

Fidelita (anglicky fidelity) vyjadřuje míru podobnosti prototypu s~finálním produktem, zejména z~hlediska vizuálního zpracování, interakčních detailů a kontextu použití.
Nízká fidelita typicky podporuje rychlé hledání směru a levné iterace, vysoká fidelita naopak umožňuje realistické testování, protože se blíží skutečné zkušenosti uživatele~\cite{joyner2022lowFidelity,flohr2023contextBased}.

LoFi prototyp v~této práci označuje ranou variantu návrhu s~omezenými vizuálními detaily, často založenou na jednoduchých rozvrženích a základní navigaci.
HiFi prototyp odpovídá pokročilejší variantě s~definovanými komponentami, typografií a chováním prvků, čímž se stává vhodným podkladem pro ověřování použitelnosti i pro navazující implementaci~\cite{joyner2022lowFidelity,flohr2023contextBased}.

\subsection{Cíl prototypu v~kontextu práce}
\label{subsec:cilProtoypu}

Cílem HiFi prototypu bylo převést výstupy uživatelské studie, scénáře případů užití a sadu požadavků do konzistentního návrhu mobilní aplikace.
Prototyp měl ověřit, zda navržené způsoby zapisování jídel, práce s~rozpoznáním a následná úprava výsledku odpovídají očekáváním uživatelů, a současně vytvořit jednoznačný podklad pro implementační část práce~\cite{cornet2020untoldStories,molina2020proposalUserCentered}.

\subsection{Rozsah a obsah HiFi prototypu}
\label{subsec:rozsahProtoypu}

HiFi prototyp byl vytvořen jako graficky plnohodnotný model mobilní aplikace a pokrývá jak klíčové toky pro zaznamenávání jídel, správu uživatelských dat a nutriční přehledy, tak i pokročilé funkce jako dotazování v~přirozeném jazyce nad daty a personalizované nastavení dietních omezení a intolerancí.
V~kapitole implementace jsou následně prezentovány konkrétní obrazovky, zatímco tato kapitola shrnuje záměr a rozsah prototypu jako výstupu návrhové fáze.

\subsection{Testování prototypu}
\label{subsec:testovaniProtoypu}

Testování prototypu je v~této práci pojato jako formativní ověřování použitelnosti, jehož cílem je odhalit problémy v~toku úloh, ve srozumitelnosti ovládání a v~porozumění výsledkům rozpoznání.
Formativní přístup je vhodný zejména v~rané fázi, protože umožňuje cíleně upravovat návrh ještě před implementací a tím snížit riziko nákladných změn v~pozdějších fázích~\cite{weichbroth2024usabilityTesting}.

Pro realizaci testování je vhodná úlohová metoda doplněná o~verbalizaci postupu typu think aloud, která poskytuje detailní vhled do mentálního modelu uživatele a do důvodů jednotlivých rozhodnutí během práce s~prototypem.
V~literatuře je think aloud běžně využíván při hodnocení mobilních zdravotních aplikací a umožňuje zachytit konkrétní bariéry v~interakci~\cite{lunde2023thinkAloud,weichbroth2024usabilityTesting}.

Testovací úlohy byly sestaveny jako zjednodušené varianty již definovaných scénářů případů užití.
Tím se eliminuje duplicitní popis funkcionalit a současně se zachovává přímá návaznost mezi analýzou potřeb, návrhem prototypu a ověřením použitelnosti~\cite{weichbroth2024usabilityTesting}.

% TODO: doplnit popisy testovacích úloh T1--T7.
\begin{description}[style=nextline]
	\item[T1] TODO
	\item[T2] TODO
	\item[T3] TODO
	\item[T4] TODO
	\item[T5] TODO
	\item[T6] TODO
	\item[T7] TODO
\end{description}

Hodnocení lze opřít o~kombinaci kvalitativních pozorování a jednoduchých kvantitativních ukazatelů, například úspěšnost dokončení úlohy, počet chyb nebo potřebu nápovědy.
Pro subjektivní hodnocení použitelnosti lze využít standardizovaný dotazník System Usability Scale, který je v~mHealth studiích běžně používán pro rychlé srovnání vnímané použitelnosti~\cite{weichbroth2024usabilityTesting,wang2022measuringMobileApp}.

Výstupem testování je seznam identifikovaných problémů a doporučení pro úpravy prototypu, ideálně seřazený podle závažnosti a dopadu na klíčové toky.
Tato struktura umožňuje transparentně propojit výsledky ověřování s~následnými designovými úpravami a s~implementačními prioritami.

\section{Shrnutí návrhu aplikace}
\label{sec:shrnutiNavrhu}

% TODO: přeformulovat, odstranit významové duplicity.

Tato práce se zaměřuje na návrh mobilní aplikace pro rozpoznávání potravin pomocí umělé inteligence a pro sledování kalorického příjmu.
Volba tématu vychází z~potřeby zjednodušit a zrychlit zapisování jídel, které je v~existujících řešeních často časově náročné a dlouhodobě obtížně udržitelné.

Zásadním vstupem byla uživatelská studie realizovaná strukturovanými rozhovory.
Získaná zpětná vazba byla pro návrh aplikace velmi přínosná, protože přinesla nové pohledy na problematiku, vyvrátila některé předpoklady a jiné naopak potvrdila.

Na základě rozhovorů byly vytvořeny scénáře případů užití, které formulují klíčové interakce uživatele se systémem a propojují kvalitativní zjištění s~návrhem aplikace.
Současně byly zpracovány funkční a nefunkční požadavky, které vymezují rozsah návrhu a cílové charakteristiky.

HiFi prototyp slouží jako konsolidovaný výstup návrhové fáze, který propojuje zjištění z~uživatelského výzkumu, formalizované scénáře případů užití a specifikované požadavky.
V~další části práce je prototyp využit jako reference pro implementaci a pro ověřování, zda výsledná aplikace naplňuje zamýšlené toky a očekávanou použitelnost.

Výstupem návrhové části je HiFi prototyp, který převádí případy užití a požadavky do konkrétní podoby uživatelského rozhraní.
Prototyp slouží jako ověřitelný podklad pro implementaci a pro testování navázané na definované případy užití.
V~další fázi práce bude cílem navázat implementací vybraných funkcí, které ověří, do jaké míry návrh naplňuje cíle definované uživatelskou studií.
